%% kapa: exam=1 round=28 year=2017 period=1 form=A subject=civil_law subject_ko=민법 seq=13 number=13
%% source: Q-Net 감정평가사 기출문제 (cst00309) artlSeq=5199603 file=extracted/1008/5199603/2017년 제28회 감정평가사 1차시험 문제지/1교시 A형.pdf
%% conversion: kordoc 4.24.51 -> tools/gen1.py (auto); review=proofread-sonnet; flags=none
\begin{question}{13}
행위능력자 甲은 대리권 없이 乙을 대리하여 乙 소유 토지를 丙에게 매도하는 계약을 체결하였다. 이에 관한 설명으로 옳은 것은? (다툼이 있으면 판례에 따름)
\begin{choices}
  \item 乙이 매매계약을 추인한 경우 다른 의사표시가 없으면 그 계약은 추인한 때부터 장래를 향하여 효력이 있다.
  \item 丙이 계약 당시 甲의 대리권 없음을 알았더라도 乙의 추인이 있을 때까지 丙은 그 계약을 철회할 수 있다.
  \item 상대방에 대한 무권대리인의 책임에 관한 규정에 의하여 甲은 丙에게 무과실책임을 진다.
  \item 丙이 계약 당시 甲의 대리권 없음을 알았다면 丙은 상당한 기간을 정하여 乙에게 추인 여부의 확답을 최고할 수 없다.
  \item 甲의 대리권 없음을 알지 못한 丙은, 乙이 甲에 대하여 매매계약을 추인한 사실을 몰랐더라도 계약을 철회할 수 없다.
\end{choices}
\end{question}
\answer{3}
